% TOP_PROBLEM: {"id": 10000037, "title": "Uniqueness of percolation on graphs roughly isometric to lattices", "category_id": 19, "category": {"id": 19, "name": "probability", "display_name": "Probability", "description": "Problems involving probability theory, stochastic processes, and random structures.", "slug": "probability", "order_index": 19, "created_at": "2026-07-31T15:26:25.671Z"}, "status": "partially_solved", "problem_number": "AMR-099-0037", "set_id": 13}
% FIRST_POSTED: 2026-10-01
% ENTRY_KIND: statement_only
% UnsolvedMath ID 10000037; source catalogue code AMR-099-0037
% !TeX program = lualatex
% Definition and sources only; this file is not an LLM solution attempt.
\documentclass[11pt,a4paper]{article}
\usepackage[margin=25mm]{geometry}
\usepackage{fontspec}
\setmainfont{lmroman10-regular.otf}[BoldFont=lmroman10-bold.otf,
 ItalicFont=lmroman10-italic.otf,BoldItalicFont=lmroman10-bolditalic.otf]
\usepackage{amsmath,amssymb,mathtools}
\usepackage{unicode-math}
\setmathfont{latinmodern-math.otf}
\AtBeginDocument{\let\setminus\smallsetminus}
\usepackage{xurl,hyperref}
\hypersetup{colorlinks=true,urlcolor=blue}
\setcounter{secnumdepth}{0}
\setlength{\parindent}{0pt}
\setlength{\parskip}{5pt}
\setlength{\emergencystretch}{3em}
\begin{document}
\section*{Uniqueness of percolation on graphs roughly isometric to lattices}\label{10000037}
{\small UnsolvedMath ID 10000037 · AMR-099-0037 · Probability · AMR Open Problem Lists}
\subsection{Definitions and mathematical statement}
Fix $d\ge1$. Let $G$ be an infinite connected graph of uniformly bounded degree which is roughly isometric to the nearest-neighbor lattice $\mathbb Z^d$. Explicitly, suppose there are constants $A\ge1$, $B\ge0$ and a map $f:V(G)\to\mathbb Z^d$ such that
\[
A^{-1}d_G(x,y)-B\le\|f(x)-f(y)\|_1\le A d_G(x,y)+B
\]
for all vertices $x,y$, and every lattice point lies within $\ell_1$-distance $B$ of some image point. This is the source's multiplicative-and-additive rough-isometry convention.

For a fixed $p\in[0,1]$, independently declare each edge of $G$ open with probability $p$. Let $N_\infty$ be the number of infinite connected components of the open-edge subgraph. Must
\[
\mathbb P_p(N_\infty\le1)=1
\]
hold for every such $G$ and every $p$?

The assertion includes parameters where no infinite component exists. When percolation does occur, it asks for uniqueness throughout the entire parameter range, including any critical value. No invariance of $G$ under a transitive group, and no translation-invariant choice of the coarse comparison map, is assumed. These omissions are central because ordinary lattice uniqueness arguments use spatial symmetries absent from an arbitrary roughly isometric graph.

The degree bound is uniform over the vertices of each graph, although it and the rough-isometry constants may depend on $G$. The problem concerns Bernoulli bond percolation with one constant parameter, rather than arbitrary dependent or spatially inhomogeneous random subgraphs.
\subsection{Short English statement}
Does every bounded-degree graph with the large-scale geometry of a Euclidean lattice have at most one infinite Bernoulli cluster?
\subsection{Sources}
\begin{itemize}
\item[S1] ULAM.AI and the UnsolvedMath Contributors, supplied problems.json, record 10000037 (AMR-099-0037), AMR Open Problem Lists. Catalogue identity and source transcription; CC BY 4.0.
\url{https://huggingface.co/datasets/ulamai/UnsolvedMath}.
\item[S2] Benjamini - Coarse Geometry and Randomness (Saint-Flour notes), Open problem 9.6, PDF page 64. Original source indexed by AMR-099-0037.
\url{https://www.wisdom.weizmann.ac.il/~itai/stflouraug24.pdf#page=64}.
\item[S3] I. Benjamini, Coarse Geometry and Randomness, primary Saint-Flour notes; the numbered question and its surrounding definitions.
\url{https://arquivo.pt/noFrame/replay/20201231041548id_/http://www.wisdom.weizmann.ac.il/~itai/stflouraug24.pdf}.
\end{itemize}
\end{document}
